Single-node IoT monitoring must distinguish environmental events from sensor faults while limiting acquisition and transmission. We evaluate TAIS, an executable integration of heuristic sensor-health diagnostics, event-aware sampling, upload suppression, stateless binary snapshots, and an independently decoded receiver. Fixed-period, delta-driven, variance-driven, and explicitly Rao-inspired schedules provide independent controls. Evaluation uses two separately specified synthetic tasks: G1 supports development and held-out comparison, whereas G2 applies frozen policies without retuning and changes signal dynamics, fault processes, event timing, and ground-truth interval definitions. SU40's F1 difference from the variance-budget control changes from +4.88 percentage points in G1 to -4.32 in G2; the G2 descriptive seed-paired 95\% interval, [-9.11, +0.18], includes zero. Thus its near-budget quality advantage is not established across these tasks. Lower application-byte counts and fewer operationally attributed fault-induced reports coexist with missed short events and older receiver data. All communication counts come from actual encoding and independent receiver reconstruction. The study contributes reproducible system integration, independent resource comparisons, and explicit failure boundaries, rather than a new trust theory or universally superior policy. Two project-defined synthetic tasks and ideal local delivery do not establish field robustness or measured battery savings.
